\documentclass[12pt,a4paper]{article}

%----------------------------------------
% PAQUETES
%----------------------------------------
\usepackage[spanish]{babel}
\usepackage[utf8]{inputenc}
\usepackage[T1]{fontenc}
\usepackage{geometry}
\usepackage{setspace}
\usepackage{titlesec}
\usepackage{fancyhdr}
\usepackage{graphicx}
\usepackage{longtable}
\usepackage{array}
\usepackage{booktabs}
\usepackage{hyperref}
\usepackage{enumitem}
\usepackage[backend=biber,style=apa]{biblatex}
\addbibresource{referencias.bib}

\geometry{margin=2.5cm}
\setstretch{1.4}

%----------------------------------------
% ENCABEZADO Y PIE
%----------------------------------------
\pagestyle{fancy}
\fancyhf{}
\fancyhead[L]{Convocatoria Nacional de Investigación Científica y Humanística 2026}
\fancyhead[R]{SECIHTI}
\fancyfoot[C]{\thepage}

%----------------------------------------
% FORMATOS DE TÍTULOS
%----------------------------------------
\titleformat{\section}{\large\bfseries}{\thesection.}{1em}{}
\titleformat{\subsection}{\normalsize\bfseries}{\thesubsection.}{1em}{}

%----------------------------------------
% PORTADA
%----------------------------------------
\begin{document}

\begin{titlepage}
    \centering
    
    % LOGO (opcional)
    % \includegraphics[width=0.3\textwidth]{logo.png}\par\vspace{1cm}
    
    {\large \textbf{SECRETARÍA DE CIENCIA, HUMANIDADES, TECNOLOGÍA E INNOVACIÓN}}\\[0.5cm]
    {\normalsize Convocatoria Nacional de Investigación Científica y Humanística 2026}\\[2cm]

    {\Huge \textbf{Prótesis mamarias externas mediante perfilometría óptica e impresión 3D}}\\[1.5cm]

    \begin{flushleft}
    \textbf{Institución beneficiaria: Universidad Autónoma de Zacateca} \\[0.2cm]
    \textbf{Responsable técnico: Dra. Susana Leticia Burnes Rudecino } \\[0.2cm]
    \textbf{Área: ingeniería, disciplina: Ingeniería Biomédica, subdisciplina: Prótesis} \\[0.2cm]
    \textbf{Fecha:} \today
    \end{flushleft}

    \vfill

\end{titlepage}

\tableofcontents
\newpage

%----------------------------------------
\section{Datos Generales del Proyecto}

\subsection{Prótesis mamarias externas mediante perfilometría óptica e impresión 3D}

\subsection{Área de conocimiento}
\begin{itemize}
    \item Área: Ingeniería
    \item Disciplina: Ingeniería Biomédica
    \item Subdisciplina: Prótesis
    \item Palabras clave: Prótesis mamarias, perfilometría óptica, impresión 3D
\end{itemize}

\subsection{Resumen ejecutivo}
La reconstrucción e impresión de prótesis mamarias externas tiene un alto potencial para el desarrollo científico, tecnológico y de servicios. En nuestro país, los servicios públicos y privados de diseño de prótesis son escasos incluso nulos en algunas regiones, siendo el factor más importante la falta de desarrollo científico y tecnológico en este campo. Este proyecto propone el desarrollo científico-tecnológico abierto para el diseño de prótesis mamarias externas en cuatro etapas; captura, edición, impresión y evaluación. En la captura se utilizarán métodos ópticos no ionizantes de alta precisión, la edición se basa en el procesamiento de nube de puntos y patrones, la impresión será directa y en molde y la evaluación se hará mediante un estudio longitudinal temporal.  La propuesta tiene cuatro puntos de investigación de frontera: (1) algoritmos para la reconstrucción. La reconstrucción 3D se ha realizado con modelos matemáticos conceptuales, estos modelos tienen un rendimiento excelente en condiciones ideales sin embargo pierden robustez en aplicaciones como la de prótesis. Se propone que los algoritmos de las etapas cruciales; detección de fase, filtrado de interferograma y calibración, sean modelos matemáticos asistidos por inteligencia artificial. (2) Relaciones antropométricas entre los senos y proporciones corporales superiores. En muchos casos la mastectomía radical se realiza en uno de los senos, por lo que se puede escanear el otro, así como el tronco superior. Aquí investigaremos la relación antropométrica que existe entre los senos y el tronco superior, se propone encontrar una función de transporte que a partir del modelo 3D del seno capturado y de la persona, se pueda reconstruir el otro. (3) Uso de materiales o biomateriales para prótesis. Se investigará materiales biocompatibles, mecánicamente adecuados, estables químicamente y durables tales como silicona médica e hidrogeles. (4) Diseño estético y ergonómico. Se trabajará con teoría estética, análisis biomecánico y postural para un diseño personalizado en la prótesis y en los sujetadores. En todas la fases proponemos desarrollo científico y tecnológico, mediante nuevo conocimiento, algoritmos, desarrollo de software, materiales y diseño.20Nuestro proyecto tiene un impacto social alto. El cáncer de mama es el segundo cáncer más frecuente y desafortunadamente la primera causa de muerte en mujeres a nivel mundial. La mastectomía radical es uno de los procedimientos utilizados y alrededor del 90\% de las mujeres que se han sometido a este tratamiento han decidido utilizar prótesis; sin embargo, estas suelen ser estandarizadas y de difícil acceso, lo que puede generar en las pacientes problemas de postura y autoestima debido a que no se adaptan adecuadamente a su cuerpo. La población beneficiada serán importante. Además que el proyecto podrá ser replicado fácilmente, incluso con algunas adecuaciones puede emplearse para otro tipos de prótesis.  

\subsection{Institución beneficiaria}
Universidad Autónoma de Zacatecas. Número de registro (Reniecyt) UAZ: 1800173
Responsable Técnica: Dra. Susana Burnes Rudecino CVU: 91003
Representante Legal: Dr. Ángel Román Gutiérrez CVU: 35415
Represponsable Administrativo: Dr. Jorge Issac Galván Tejada CVU: 326164

\subsection{Persona responsable técnica}
\begin{itemize}
    \item Nombre completo:Susana Burnes Rudecino 
    \item Grado académico: Doctorado
    \item Perfil Rizoma: CVU: 91003
    \item Información demográfica:
    \item Adscripción: Universidad Autónoma de Zacatecas
\end{itemize}

\subsection{Participantes}
\begin{longtable}{p{4cm} p{3cm} p{4cm} p{3cm}}
\toprule
Nombre & Grado & Perfil Rizoma & Institución \\
\midrule
Nombre 1 & Grado & Perfil & Institución \\
Nombre 2 & Grado & Perfil & Institución \\
\bottomrule
\end{longtable}

%----------------------------------------
\section{Introducción, objetivos, metodología y plan de trabajo}

\subsection{Antecedentes}
El uso de las prótesis mamarias se remonta a la antigüedad; sin embargo, su evolución ha estado marcada por avances significativos en materiales, diseño y enfoque clínico. Históricamente, desde el siglo XVI se empleaban rellenos de tela, algodón y fibras vegetales para simular el volumen mamario tras la pérdida de tejido. Posteriormente, durante el siglo XX, especialmente tras la Primera Guerra Mundial, se introdujeron materiales como caucho y tejidos naturales, representando un avance en la funcionalidad de estos dispositivos \cite{ACS2023}.
En la actualidad, el desarrollo de prótesis mamarias externas se sitúa en la intersección de la biomecánica, la ingeniería biomédica y la rehabilitación física. En los últimos cinco años, se observa una clara tendencia hacia la personalización, la optimización del confort y la reducción de efectos secundarios asociados al uso prolongado, como dolor musculoesquelético y alteraciones posturales \cite{WHO2021}

El cáncer de mama es una de las principales causas que conducen al uso de estas prótesis. Para combatir esta enfermedad, los tratamientos incluyen quimioterapia, radioterapia, terapia hormonal y/o mastectomía. La mastectomía consiste en la extirpación del tejido mamario afectado, lo que altera la apariencia física y puede impactar negativamente en la autoestima de las mujeres, generando ansiedad y depresión, debido a que el seno es un símbolo de feminidad y maternidad \cite{Jetha2017} falta cita kang2024cross.

A pesar de la existencia de procedimientos de reconstrucción mamaria, pocas mujeres acceden a ellos debido a su alto costo, salvo en casos donde se cuenta con apoyo gubernamental o financiamiento externo faltan citas:{lai2024robotic,novo2014coste}. Por ello, tras el impacto emocional posterior a la mastectomía, muchas pacientes optan por el uso de prótesis mamarias externas falta cita: {wiedemann2017external}.
En las etapas iniciales posteriores a la cirugía, se emplean prótesis ligeras de espuma o relleno textil. Posteriormente, entre seis y ocho semanas después, una vez cicatrizada la herida, se utilizan prótesis de silicona o elaboradas de manera personalizada con un peso adecuado falta cita: {crompvoets2012prosthetic}. Estas prótesis actúan como sustitutos del seno natural, contribuyendo a mejorar la salud física y psicológica, incluyendo la autoimagen, la autoestima y la percepción de feminidad.

La correcta selección de la prótesis es fundamental, ya que existen múltiples opciones en cuanto a tamaño, forma, peso y material. Se ha demostrado que el uso prolongado de prótesis mamarias externas incrementa la satisfacción emocional y mejora la calidad de vida en mujeres que no se han sometido a reconstrucción mamaria, además de reducir el estrés asociado a la mastectomía faltann citas: {glaus2009long,hopwood2001body}. Asimismo, con el tiempo, el uso de la prótesis favorece la autoaceptación tras la cirugía falta cita: (Nissen, 2001).

El progreso en el diseño de prótesis ha permitido el desarrollo de una amplia variedad de opciones: con o sin pezón, con diferentes texturas, colores, sistemas autoadhesivos e incluso modelos especializados para actividades deportivas faltan citas: {fitch2012perspectives,livingston2005women,liang2015factors}. No obstante, los altos costos continúan siendo una barrera importante, lo que limita el acceso a prótesis adecuadas faltan citas:{healey2003external,krishna2022comments}. Además, su vida útil —aproximadamente dos años en el caso de las prótesis de silicona— representa un gasto recurrente, al igual que el uso de sostenes especializados necesarios para su colocación faltan citas:{gallagher2010external,mcghee2020effect}.

En contextos locales como el estado de Zacatecas, México, muchas mujeres adquieren prótesis a través de instituciones que las proporcionan mediante campañas de apoyo. Sin embargo, estas suelen ser estandarizadas, con formas y pesos que no siempre se adaptan adecuadamente a las características individuales, afectando la comodidad, la naturalidad y la movilidad.

Diversos estudios recientes han señalado que las prótesis tradicionales, especialmente aquellas con mayor peso, pueden provocar compensaciones posturales, sobrecarga en la región cervical y dorsal, así como alteraciones en la marcha \cite{Leung2017}. En respuesta, la investigación actual se ha enfocado en el desarrollo de prótesis más ligeras mediante siliconas de baja densidad, geles avanzados y estructuras internas diseñadas para una mejor redistribución del peso.

Una de las innovaciones más relevantes es la incorporación de tecnologías de digitalización, como el escaneo tridimensional y la impresión 3D, que permiten fabricar prótesis altamente personalizadas. Estas tecnologías facilitan la reproducción precisa de la anatomía de la paciente, mejorando la simetría y la adaptación al contorno corporal, lo que impacta positivamente en la postura y en la percepción de la imagen corporal \cite{Singh2020,Tan2023}.
En términos de diseño, las prótesis contemporáneas integran materiales hipoalergénicos, superficies transpirables y sistemas de adhesión mejorados que permiten su fijación directa a la piel, reduciendo la dependencia de prendas especializadas \cite{WHO2021}. Asimismo, se han desarrollado prótesis parciales y modulares que permiten ajustes progresivos según las necesidades de cada paciente.

Finalmente, el estado del arte refleja un cambio importante en el enfoque clínico: ya no se prioriza únicamente la estética, sino un abordaje integral centrado en la calidad de vida. Esto incluye la evaluación de la postura, el equilibrio corporal, la prevención del dolor y el acompañamiento en la adaptación física y emocional tras la mastectomía \cite{McNeely2021,Hojan2016}.

En conjunto, estos avances evidencian una transición hacia prótesis mamarias más funcionales, personalizadas y orientadas al bienestar integral, integrando principios biomecánicos y tecnológicos para mejorar la experiencia de uso y reducir las complicaciones asociadas.

\subsection{Pertinencia}
 El uso de prótesis mamarias después de un tratamiento de cáncer de mama donde el tratamiento requirió la extirpación total del seno, produce fuertes efectos positivos en las mujeres que llevaron el proceso de cáncer: como mejorar la autoestima, corrección de la asimetría mamaria, así como por razones estéticas.

El desarrollar prótesis mamarias personalizadas para un grupo de mujeres de Zacatecas que tienen mastectomía radical, permitirá que tengan fácil acceso a prótesis cómodas, acorde a su anatomía específica, permitiendo que tengan una mejor salud física y emocional. Pues las prótesis al venir en tallas no corresponden en peso y forma para cada paciente

\subsection{Hipótesis o preguntas de investigación}
Planteamiento claro.

\subsection{Objetivo general}
Obtener una prótesis personalizada de bajo costo para las mujeres zacatecanas.

\subsection{Objetivos específicos}
\begin{enumerate}
    \item Hacer un estudio de campo sobre las afectaciones de usar prótesis de mama en un grupo de UNEME Zacatecas.
    \item Obtener el valor aproximado de una mama para llevar a cabo una prótesis personalizada de bajo costo usando el volumen y peso de la mama de cada paciente.
    \item Realizar el escaneo de la mama femenina para llevar a cabo una prótesis personalizada de bajo costo
\end{enumerate}

\subsection{Metas}
\begin{longtable}{p{3cm} p{7cm} p{3cm}}
\toprule
Etapa & Meta & Indicador \\
\midrule
1 & Descripción & Métrica \\
2 & Descripción & Métrica \\
\bottomrule
\end{longtable}

\subsection{Metodología}
Descripción detallada de métodos, técnicas y modelos.

\subsection{Cronograma de actividades}
\begin{longtable}{p{2cm} p{6cm} p{3cm} p{2cm}}
\toprule
Etapa & Actividad & Responsable & Tiempo \\
\midrule
1 & Actividad & Nombre & Meses \\
2 & Actividad & Nombre & Meses \\
\bottomrule
\end{longtable}

\subsection{Factores de riesgo y estrategias}
\begin{itemize}
    \item Riesgo: Mitigación
    \item Riesgo: Mitigación
\end{itemize}

\subsection{Bibliografía}
\begin{thebibliography}{9}
\bibitem{ref1} Autor, título, revista, año.
\end{thebibliography}

%----------------------------------------
\section{Resultados esperados e impacto}

\subsection{Resultados esperados}
Descripción.

\subsection{Impacto social}
Impacto potencial.

%----------------------------------------
\section{Productos entregables comprometidos}

\subsection{Académicos}
Artículo en revista mexicana de acceso abierto.

\subsection{Acceso universal del conocimiento}
Divulgación científica.

\subsection{Formación de capacidades}
Formación de estudiantes, cursos, mentoría.

%----------------------------------------
\section{Presupuesto detallado}

\subsection{Desglose}
\begin{longtable}{p{2cm} p{5cm} p{3cm} p{2cm}}
\toprule
Etapa & Concepto & Justificación & Monto \\
\midrule
1 & Concepto & Justificación & \$ \\
2 & Concepto & Justificación & \$ \\
\bottomrule
\end{longtable}

\subsection{Justificación}
Explicación del uso de recursos.

%----------------------------------------
\section{Propuesta de evaluadores}

\begin{itemize}
    \item Nombre -- Institución -- Área
    \item Nombre -- Institución -- Área
    \item Nombre -- Institución -- Área
\end{itemize}

%----------------------------------------
\section{Carta Oficial de Postulación}

Se deberá incluir carta firmada electrónicamente que contemple:

\begin{itemize}
    \item Apoyo institucional
    \item Congruencia institucional
    \item Designación de responsables
    \item Acuerdos de colaboración
    \item Cumplimiento normativo
    \item Declaración de no adeudos
\end{itemize}

\printbibliography
%----------------------------------------
\section{Documentos Anexos}

\subsection{Protocolo de investigación}
Se pueden incluir figuras, gráficas, fórmulas, etc.

%----------------------------------------

\end{document}